\chapter{Metodologia}
\label{chap:metodologia}

\section{Delineamento de pesquisa}
Com base nos objetivos da presente pesquisa, a mesma se classifica como explicativa. Esse tipo de pesquisa "têm como preocupação central identificar os fatores que determinam ou que contribuem para a ocorrência dos fenômenos." \citem[p. 87]{Gil}. Ao analisarmos o desempenho dos diferentes modelos econométricos voltados para a predição do PIB brasileiro, estamos buscando identificar e compreender os diferentes fatores que explicam esse desempenho, enquadrando-a como explicativa.

Sobre os procedimentos de coleta de dados, a presente pesquisa se enquadra como documental, por utilizar-se de dados que não receberam um tratamento e estão dispostos em repositórios digitais de dados. Segundo \citeonline{Gil} a pesquisa documental apropria-se de diversas fontes, não somente material impresso em bibliotecas, onde os documentos utilizados podem ter um tratamento analítico ou não.

Por fim, ao se tratar de uma pesquisa que compara diferentes modelos econométricos, assim como avalia o desempenho dos mesmos, seguindo o rigor matemático e estatístico, ela se enquadra como uma pesquisa predominantemente quantitativa onde apropria-se do método comparativo e estatístico. Para \citeonline{LAKATOSMARCONI} o método estatístico trata-se da redução dos fenômenos a termos quantitativos e a manipulação estatística, onde possibilita a generalizações sobre sua natureza, ocorrência ou significado. Já "o método comparativo é usado tanto para comparações de grupos." \citem[p. 107, 108]{LAKATOSMARCONI}.


%	Aqui eu tenho que falar sobre os objetivos
%	O procedimento de coleta
%	Quanto a abordagem as questões de Pesquisa
\section{Procedimento de coleta de dados}

As variáveis foram escolhidas baseadas na dissertação de \citeonline{rebelo2022interpretabilidade}, descritas mais evidentemente na sessão \ref{cap:revisao-literatura}. As variáveis escolhidas podem ser verificadas na tabela \ref{tab:fontes_dados}. Os dados utilizados na pesquisa foram extraídos de duas APIs (\textit{Application Programming Interface}, no português, Interface de programação de aplicações). A relação variável - Fonte - URL da API pode ser vericada na mesma tabela.

A primeira fonte foi a API de dados agregados do IBGE, onde nela existem informações econômicas, sociais e demográficas do Brasil. 


\begin{citacao}
	O IBGE possui uma Application Program Interface (API) de serviços de dados abertos que disponibiliza dados processáveis por máquina para qualquer cidadão ou sistema, permitindo a interoperabilidade e o enriquecimento a partir dos dados abertos disponibilizados. \citem[p. 2]{silva2024dados}.
\end{citacao}

A segunda API utilizada foi a API do BACEN (Banco Central do Brasil), onde foram extraídos dados financeiros. Segundo o próprio \citeonline{bcb_dados_abertos} o "catálogo de dados, permite que os usuários encontrem os conjuntos de informações de seu interesse, entendam sua estrutura e suas limitações, e tenham acesso ao caminho para os dados."

O período analisado foi de 2007 à 2022, devido a disponibilidade de datas nas bases de dados. 

\begin{table}
	\centering
	\caption{Fontes de dados utilizadas nas APIs do IBGE e BACEN}
	\label{tab:fontes_dados}
	\resizebox{\textwidth}{!}{
		\begin{tabular}{p{4cm} p{2cm} p{9cm}}
			\hline
			\textbf{Indicador} & \textbf{Fonte} & \textbf{URL / Endpoint da API} \\
			\hline
			PIB & IBGE & \small \url{https://servicodados.ibge.gov.br/api/v3/agregados/6784/periodos/-20/variaveis/9808?localidades=N1[all]} \\
			Deflator do PIB & IBGE & \small \url{https://servicodados.ibge.gov.br/api/v3/agregados/6784/periodos/-20/variaveis/9811?localidades=N1[all]} \\
			população (milhões) & IBGE & \small \url{Ainda nao coletado} \\
			índice de direitos de propriedade & IBGE & \small \url{Ainda nao coletado} \\
			índice de eficiência jurídica & IBGE & \small \url{Ainda nao coletado} \\
			índice de integridade governamental & IBGE & \small \url{Ainda nao coletado} \\
			índice de despesa governamental  & IBGE & \small \url{Ainda nao coletado} \\
			despesa pública em educação (\% PIB) & IBGE & \small \url{Ainda nao coletado} \\
			despesa pública em saúde (\% PIB) & IBGE & \small \url{Ainda nao coletado} \\
			despesa pública em proteção ambiental (\% PIB) & IBGE & \small \url{Ainda nao coletado} \\
			índice de carga fiscal & IBGE & \small \url{Ainda nao coletado} \\
			índice de saúde fiscal  & IBGE & \small \url{Ainda nao coletado} \\
			dívida publica (\% PIB) & IBGE & \small \url{Ainda nao coletado} \\
			índice de liberdade empresarial & IBGE & \small \url{Ainda nao coletado} \\
			índice de liberdade laboral & IBGE & \small \url{Ainda nao coletado} \\
			taxa de desemprego & IBGE & \small \url{Ainda nao coletado} \\
			índice de liberdade monetária  & IBGE & \small \url{Ainda nao coletado} \\
			taxa de inflação  & IBGE & \small \url{Ainda nao coletado} \\
			índice de liberdade comercial & IBGE & \small \url{Ainda nao coletado} \\
			balança corrente  & IBGE & \small \url{Ainda nao coletado} \\
			Taxa aduaneira & IBGE & \small \url{Ainda nao coletado} \\
			índice de liberdade investimento  & IBGE & \small \url{Ainda nao coletado} \\
			investimento direto estrangeiro (\% PIB) & IBGE & \small \url{Ainda nao coletado} \\
			índice de liberdade financeira & IBGE & \small \url{Ainda nao coletado} \\
			\hline
		\end{tabular}
	}
\end{table}




\section{Procedimento de análise de dados}

A análise dos dados é um processo muito importante ao usar modelos matemáticos, conforme \citeonline{MoniqueGIGO} a qualidade das respostas depende diretamente dos dados usados 
\footnote{O termo coloquial em inglês para essa situação se chama GIGO ou \textit{Garbage In Garbage Out}, que em tradução livre - lixo entra lixo sai.}.
Ou seja, não basta ter fontes confiáveis como o \citem{bcb_dados_abertos} e o \citem{ibge_api}, é necessário um procedimento confiável de ponta a ponta.
Para esse processamento dos dados, foi utilizado a linguagem de programação \textit{Python}
\footnote{Python é amplamente utilizado em computação científica e numérica \citem{python}}
e suas bibliotecas.


O primeiro processo de tratamento foi a conversão dos tipos nativos das APIs para um formato adequado para análise com Python, isso foi feito através da biblioteca Pandas. "Pandas é uma biblioteca de Python de estruturas de dados e ferramentas estatísticas, inicialmente desenvolvida para aplicações de finanças quantitativas." \citem[p. 1, \textit{tradução nossa}]{mckinney2010pandas}.

O segundo passo de tratamento foi a inflação dos dados para 2022, último ano da série, isso foi feito através de um índice de preços, que pode verificado na equação \ref{eq:inflacao}. O \citeonline{ibge_pib_series} traz a definição de deflator do PIB, porém que se aplica as demais variáveis, em que um deflator é a variação do valor corrente da variável no tempo $(t)$ em relação ao valor da varável do ano $(t)$ a preços do ano anterior, ou seja a variação percentual em relação ao mesmo período do ano anterior.
\begin{equation}
	\label{eq:inflacao}
	\text{Valor real} = \frac{\text{Valor Nominal}}{\text{índice de Preços}} \times 100
\end{equation}

O terceiro passo de tratamento foi o treinamento dos modelos. Nesse passo foram utilizados as bibliotecas Scikit-learn e Tensorflow, os modelos foram treinados com dados de 2007 à 2017, totalizando 10 anos na série histórica (dados de treinamento), e deixando os últimos 5 anos (2017 à 2022) para a testagem dos resultados do modelo. Para \citeonline{santos2022previsao} as eficácias dos modelos nos dados de treinamento não são importantes, o que é relevante é como o modelo se sai com os dados de teste, ou seja, se o modelo consegue ser extensivo no período fora da amostra.

O último passo foi o cálculo das métricas de avaliação dos modelos, para esse passo foram avaliadas três dimensões: acurácia, interpretabilidade e custo computacional.

Para a acurácia foram escolhidas as métricas que foram utilizadas por \citeonline{rebelo2022interpretabilidade}: o  Erro Quadrático Médio (EQM) descrito na equação \ref{eq:EQM}; Erro Absoluto Médio (EAM) descrito na equação \ref{eq:EAM}; a Raiz do Erro Quadrático Médio (REQM) descrito na equação \ref{eq:REQM}; e por fim, o erro percentual absoluto médio (EPAM) descrito na equação \ref{eq:EPAM}.

\begin{equation}
	\label{eq:EQM}
	\text{EQM} = \frac{1}{n}\sum_{i=1}^{n}(y_i \hat{y_i})^2
\end{equation}

\begin{equation}
	\label{eq:EAM}
	\text{EQM} = \frac{1}{n}\sum_{i=1}^{n}|y_i \hat{y_i}|
\end{equation}

\begin{equation}
	\label{eq:REQM}
	\text{REQM} = \sqrt{\text{EQM}}
\end{equation}

\begin{equation}
	\label{eq:EPAM}
	\text{EPAM} = \frac{1}{n}\sum_{i=1}^{n}|\frac{y_i \hat{y_i}}{y_i}| \times 100
\end{equation}

Para \citeonline{kava2022alem} a interpretabilidade está diretamente relacionada a facilidade que um ser humano tem de entender os resultados de um modelo. Para essa dimensão foram usadas as métricas: Efeitos Locais Acumulados (ELA) descrito na equação \ref{eq:ELA}; e a interação ente preditores através da estatística-H, descrita na equação \ref{eq:H}, ambas descritas por \citeonline{kava2022alem}.

\begin{equation}
	\label{eq:ELA}
	\tilde{f}_j,\text{ELA}(x) = \sum_{k=1}^{k_j(x)}\frac{1}{n_k(k)} \sum_{i:X_j \in N_j(k)}[\hat{f}(z_{k,j},x_j^{(i)})-\hat{f}(z_{k-1,j},x_j^{(i)})]
\end{equation}



\begin{equation}
	\label{eq:H}
	H_j^2 = \frac{ \sum_{i=1}^{n}[\hat{f}(x^{(i)})-PD_J(X_J^{(I)})-PD_{-J}(X_{-J}^{(I)})]^2}{\sum_{i=1}^{n}\hat{f^2}(x^{(i)})}
\end{equation}

A última dimensão é o custo computacional, essa se baseará principalmente no tempo de execução dos modelos, e na contagem de operações com ponto flutuantes - FLOPs (do inglês, FLOating Point Operations) em ambas as métricas, quanto mais alto o valor, maior é o custo computacional.


\section{Constructo da pesquisa}
A estrutura do constructo pode ser verificado na tabela \ref{tab:Constructo}
%	\begin{table}
	%		\centering
	%		\caption{Constructo da pesquisa}
	%		\label{tab:Constructo}
	%		\begin{tabular}{p{2cm} p{4cm} p{9cm}}
		%			\hline
		%			\textbf{Objetivo} & \textbf{Onde analisar} & \textbf{Como analisar}\\
		%			a & Revisão de literatura & Levantar na literatura os principais modelos econométricos tradicionais (VAR, ARIMA, DFM e MIDAS) e de IA aplicados à predição do PIB \\ 
		%			a & Revisão de literatura & Descrever cada modelo identificado, apresentando suas formulações algébricas e funcionais \\ 
		%			b & Revisão de literatura & Identificar na literatura as variáveis determinantes na estimativa do PIB \\ 
		%			b & Metodologia & Selecionar e justificar as variáveis que serão utilizadas, bem como suas respectivas fontes de dados \\ 
		%			c & Análise e Discussão dos resultados & Aplicar os modelos econométricos e de IA utilizando as variáveis selecionadas para estimar o PIB \\ 
		%			c & Análise e Discussão dos resultados & Calcular as métricas de desempenho (acurácia, interpretabilidade e custo computacional) de cada modelo \\ 
		%			d & Análise e Discussão dos resultados & Comparar os resultados obtidos entre os modelos de IA e os modelos econométricos tradicionais \\ 
		%			d & Conclusão & Discutir criticamente o desempenho dos modelos e destacar suas vantagens e limitações em cada dimensão de avaliação \\ 
		%			\hline
		%		\end{tabular}
	%	\end{table}
\begin{table}[H]
	\centering
	\caption{Constructo da pesquisa}
	\label{tab:Constructo}
	\begin{tabular}{p{14.5cm}}
		\hline
		\textbf{Pergunta de Pesquisa} \\
		\hline
		Em que medida modelos preditivos baseados em \ac{ia}, comparados aos modelos tradicionais VAR, ARIMA, DFM e MIDAS, conseguem estimar o PIB brasileiro? \\
		\hline
		\textbf{Objetivo Geral} \\
		\hline
		Analisar em que medida modelos preditivos baseados em \ac{ia}, comparados aos modelos tradicionais VAR, ARIMA, DFM e MIDAS, conseguem estimar o PIB brasileiro.
	\end{tabular}
	\resizebox{\textwidth}{!}{
		\begin{tabular}{p{4cm} p{5.0cm} p{3.5cm} p{4.6cm}}
			\hline
			\textbf{Objetivos específicos} & \textbf{variáveis/dados à coletar} & \textbf{Como analisar} & \textbf{Fonte}. \\
			\hline
			Caracterizar os modelos econométricos tradicionais (VAR, ARIMA, DFM e MIDAS) e de IA para a predição do PIB. & Modelos econométrico com capacidade de predizer séries temporais. & Buscar através da literatura modelos já utilizados com esses fins. & \citeonline{castelli2022modelo}, \citeonline{azevedo2020modelagem}, \citeonline{azevedo2020modelagem}, \citeonline{ciquini2023modelo}, \citeonline{silva2024dados}, \citeonline{alves2009nucleo}, \citeonline{moura2020fragilidade} e \citeonline{micheletti2024avaliaccao}. \\
			\hline
			Identificar as diferentes variáveis determinantes na estimativa do PIB. & Variáveis com possível influência no PIB com facilidade de coleta e atualização recorrente. & Procurar na literatura quais variáveis são utilizadas para a predição do PIB. & \citeonline{rebelo2022interpretabilidade}, \citeonline{ibge_api} e \citeonline{bcb_dados_abertos}. \\
			\hline
			Aplicar diferentes modelos econométricos tradicionais e de IA para a predição do PIB. & Modelos econométricos e variáveis encontradas nos primeiros dois objetivos. & treinar os diferentes modelos usando os dados coletados através da linguagem de programação Python e suas bibliotecas de estatística.  & \citeonline{castelli2022modelo}, \citeonline{azevedo2020modelagem}, \citeonline{azevedo2020modelagem}, \citeonline{ciquini2023modelo}, \citeonline{silva2024dados}, \citeonline{alves2009nucleo}, \citeonline{moura2020fragilidade} e \citeonline{micheletti2024avaliaccao}.\\
			\hline
			Comparar o desempenho dos modelos de IA e econométricos tradicionais usando diferentes métricas, como desempenho computacional, precisão e interpretabilidade. & erro Quadrático Médio (EQM), Erro Absoluto Médio (EAM), a Raiz do Erro Quadrático Médio (REQM), erro percentual absoluto médio (EPAM), Efeitos Locais Acumulados (ELA), estatística-H, tempo de execução e FLOPs  & Submeter os modelos a testes nas métricas de acurácia, interpretabilidade e custo computacional. & \citeonline{rebelo2022interpretabilidade} e \citeonline{kava2022alem}.\\
			\hline 
		\end{tabular}
	}
\end{table}